\begin{description}
\item[(a)] As ``Y'' is smaller, yet hotter, it must be on the main sequence.
  \hfill(1.0)
\item[(b)]
  \begin{description}
  \item[(I)] For this we have to measure the total duration of eclipse and
    compare it with the bottom flat part (total eclipse phase) in the
    lightcurve. Total eclipse time should be $(1.25 \pm 0.10)$ days, flat
    period of eclipse $(0.25 \pm 0.05)$ days. \hfill(1.0)
    \[ \frac{r_X + r_Y}{r_X - r_Y} = \frac{1.25 \pm 0.10}{0.25 \pm 0.05}
       = 5.0 \pm 1.5 \]
    \[ \therefore\ \frac{r_X}{r_Y} = \frac{5 + 1}{5 - 1} = 1.5 \pm 0.2 \]
    \hfill(1.0)
  \item[(II)] Because star Y is hotter, star Y will be brighter than a part
    of star X with the same area. Thus, the deeper dip in the light curve is
    due to star Y passing behind star X. Thus, the flux from star X must be
    $10 \times 10^{-12}\ \mathrm{W/m^2}$. As the combined flux is
    $16 \times 10^{-12}\ \mathrm{W/m^2}$, the flux of star Y is
    $6 \times 10^{-12}\ \mathrm{W/m^2}$. Thus
    \[ \frac{L_X}{L_Y} = \frac{10 \times 10^{-12}}{6 \times 10^{-12}}
       = 1.67 \]
    \hfill(2.0)
  \end{description}
\item[(c)]
  \begin{itemize}
  \item Light Curve B: We see that the dips in brightness both become
    shallower and have become rounded at the bottom, while the maximum
    brightness is still $16 \times 10^{-12}\ \mathrm{W/m^2}$. If the radius
    of one of the stars had changed (rounded bottom), then we would expect
    the total eclipse period to also change. But that has not happened.
    Brightening of star X can explain a shallower dip but not a rounded
    trough. This means that the stars no longer fully eclipse each other.
    Hence,

    (xi): The inclination of the system relative to the Earth has changed.
    \hfill(4.0)
  \item Light Curve C: The y-axis shows that the overall flux has increased.
    Since the flux of both the stars has gone up by the same factor and there
    is no substantial change in the profile of eclipse,

    (xii): The distance of the system from the Earth has decreased.
    \hfill(4.0)
  \item Light Curve D: The regular fluctuations in brightness indicate that
    one of the components of the system is a variable star. Since the
    fluctuations persist through all the dips in flux, this indicates that
    the variable star is never completely eclipsed. Therefore,

    (ix): Star X is a variable star. \hfill(3.0)
  \item Light Curve E: The overall flux and dips in the flux are the same
    size, but occur much more frequently. This means that

    (xv): The orbital period of the system decreased. \hfill(2.0)
  \item Light Curve F: The regular fluctuations again indicate that one of
    the components of the system is a variable star. During the primary
    eclipse, the light curve is perfectly flat. This can only occur if,

    (x): Star Y is a variable star. \hfill(2.0)
  \end{itemize}
\end{description}
